\chapter{Da gravação das
reuniões}

Todas as reuniões do CRP~SP devem ser gravadas para fins de registro
posterior, não sendo permitido o compartilhamento das gravações com
terceiros ou com as partes sem autorização expressa de todos os
participantes.

Quanto à proteção dessas gravações, fica determinado que:

\begin{enumerate}

\item
    somente a Unidade de Secretaria terá acesso às gravações para registro
    das reuniões, sendo obrigatória a exclusão do arquivo após a aprovação
    do documento;
\item
    em caso de ata ou registro aprovado ao final da própria reunião, a
    gravação deve ser imediatamente excluída;
\item
    em casos de gravação para fins de treinamento, capacitação e demais
    ações pedagógicas, todas/os as/os participantes devem autorizar
    expressamente sua divulgação no início da atividade, em conjunto com o
    registro da finalidade e da limitação de tal compartilhamento;
\item
    os assuntos discutidos nas reuniões ordinárias, extraordinárias,
    assembleias ou similares e os registros de suas informações são
    passíveis de divulgação e de acesso, com exceção das partes que
    expressamente tiverem sigilo previsto em hipóteses legais e/ou possam
    violar os princípios da Lei Geral de Proteção de Dados Pessoais (LGPD;
    Lei nº~13.709/2018).
\end{enumerate}